% © 2025 Ing. David Jaroš — CC BY-NC-ND 4.0
%
% This work is licensed under a Creative Commons Attribution-NonCommercial-NoDerivatives
% 4.0 International License (CC BY-NC-ND 4.0).
%
% License History: Earlier drafts (up to v0.3) were released under CC BY 4.0.
% From v0.4 onward, all material is released under CC BY-NC-ND 4.0 to protect
% the integrity of the theoretical work during ongoing academic development.

\documentclass[12pt]{article}
\usepackage{amsmath,amssymb,amsthm}
\usepackage{geometry}
\usepackage{hyperref}
\usepackage{tcolorbox}
\usepackage{xcolor}
\geometry{margin=1in}

% Theorem-like environments
\newtheorem{conjecture}{Conjecture}
\newtheorem{proposition}{Proposition}
\newtheorem{definition}{Definition}
\newtheorem{remark}{Remark}

\title{ST-4: Chronology Protection in\\
Unified Biquaternion Theory\\[0.5em]
\large\textsc{Speculative Extension}}
\author{David Jaroš}
\date{2026-03-01}

\begin{document}
\maketitle

\begin{tcolorbox}[colback=yellow!10, title=Status]
\textbf{SPECULATIVE EXTENSION} — This content is not part of
the canonical UBT framework. It is an exploratory investigation
and has not been validated against the canonical field equations.
\end{tcolorbox}

\vspace{1em}
\noindent\textit{Output of Sub-task ST-4 from the CTC/Causal-Structure investigation track.
See also: \texttt{complex\_time\_causal\_structure.tex},
\texttt{rotation\_complex\_time.tex},
\texttt{ctc\_conditions.tex}, \texttt{ctc\_methodology.tex}.}

%----------------------------------------------------------------------
\section*{Abstract}
%----------------------------------------------------------------------
Hawking's chronology protection conjecture \cite{Hawking1992} relies on
QFT assumptions that may not hold in Unified Biquaternion Theory (UBT)
with complex time $\tau = t + i\psi$.  We examine whether (a) the UBT
imaginary sector $\psi$ acts as a natural regulator that enforces
chronology protection more strongly, or (b) it circumvents the
standard QFT argument.  All results are stated as conjectures with
explicit assumptions; no claim constitutes a proof.
All conclusions are \textbf{speculative}.

%----------------------------------------------------------------------
\section{Hawking's Chronology Protection Conjecture (Standard GR)}
%----------------------------------------------------------------------
Hawking \cite{Hawking1992} proposed that the laws of physics prevent
the formation of closed timelike curves.  The central physical mechanism
is the divergence of the renormalised stress-energy tensor
$\langle T_{\mu\nu}\rangle_{\mathrm{ren}}$ on the Cauchy horizon that
would be needed to create a CTC.  As the horizon forms, quantum fields
in the background geometry contribute an infinite back-reaction that
destroys the Cauchy horizon before CTCs can form.

\paragraph{Key QFT assumptions of the standard conjecture.}
\begin{enumerate}
  \item Spacetime is globally hyperbolic before the CTC formation epoch.
  \item Quantum fields propagate on a fixed classical background metric
    $g_{\mu\nu}$ (semiclassical gravity).
  \item The renormalised vacuum stress-energy is well-defined and its
    divergence on the compactly generated Cauchy horizon signals a
    breakdown of the semiclassical approximation.
  \item The energy conditions (or their quantum equivalents) hold in the
    bulk away from the horizon.
\end{enumerate}

%----------------------------------------------------------------------
\section{Do These Assumptions Hold in UBT?}
%----------------------------------------------------------------------
\subsection{Assumption 1: Global Hyperbolicity}
In UBT the spacetime manifold carries complex time $\tau = t + i\psi$.
If $\psi$ is treated as an internal phase (Interpretation B of
\texttt{complex\_time\_causal\_structure.tex}), then the real
$3+1$-dimensional slice of spacetime inherits the standard causal
structure and global hyperbolicity is preserved.

Under Interpretation A ($\psi$ as second time direction), the spacetime
acquires an effectively higher-dimensional structure and global
hyperbolicity in the standard sense no longer applies.

\textit{Verdict:} Under Interpretation B, Assumption 1 holds.  Under
Interpretation A, it does not.

\subsection{Assumption 2: Semiclassical Gravity}
UBT is an attempt to unify GR and QFT at the level of the field
$\Theta$.  In UBT, the metric is not a fixed classical background but
a dynamical biquaternionic object $\mathcal{G}_{\mu\nu}$.  The
semiclassical approximation — quantum fields on a classical background —
is replaced by a fully coupled biquaternionic field equation.

\textit{Verdict:} Assumption 2 does \emph{not} hold directly in UBT.
The standard QFT-on-curved-spacetime framework must be generalised to
QFT on a biquaternionic background before the divergence argument
can be evaluated.

\subsection{Assumption 3: Divergence of Renormalised Stress-Energy}
In standard QFT the divergence of $\langle T_{\mu\nu}\rangle_{\mathrm{ren}}$
on the Cauchy horizon arises from high-frequency modes of quantum fields
propagating on the CTC-enabling background.  In UBT, the imaginary
sector $\psi$ introduces a natural scale $R_\psi$ (the compactification
radius of $\psi$).  High-frequency modes with wave-vector components
$k_\psi \gg R_\psi^{-1}$ are effectively regulated by the compact $\psi$
topology.

\textit{Verdict:} Assumption 3 may be modified in UBT: the $\psi$-sector
provides a UV regulator for modes in the imaginary-time direction, which
could soften the divergence of $\langle T_{\mu\nu}\rangle_{\mathrm{ren}}$.

\subsection{Assumption 4: Energy Conditions in the Bulk}
As established in \texttt{ctc\_conditions.tex}
(Conjecture~2), the imaginary sector $h_{\mu\nu}$ can contribute
negative effective energy density to $T_{\mu\nu}^{\mathrm{eff}}$.  This
means the weak energy condition (WEC) can be violated in UBT without
introducing exotic matter by hand.

\textit{Verdict:} Assumption 4 does \emph{not} automatically hold in UBT.
Local WEC violation is possible through the imaginary metric sector.

%----------------------------------------------------------------------
\section{Possible Outcomes for Chronology Protection in UBT}
%----------------------------------------------------------------------
Given the analysis above, two scenarios are possible:

\subsection{Scenario A: Enhanced Chronology Protection}
If the imaginary sector $\psi$ provides a UV regulator (Assumption 3
modified) and the energy conditions hold on average (Assumption 4
replaced by an averaged energy condition), then the standard divergence
argument might be \emph{strengthened} in UBT: the $\psi$-regulator
smears out the Cauchy-horizon divergence in the imaginary direction but
does not remove it, and back-reaction still destroys the horizon before
CTCs form.

\begin{conjecture}[Enhanced chronology protection]
\label{conj:enhanced_cp}
\textbf{Assumptions:}
(a) $\psi$ is an internal phase coordinate (Interpretation B);
(b) the $\psi$-compactification radius $R_\psi$ is finite;
(c) the averaged null energy condition holds in the real sector;
(d) UBT reduces to semiclassical GR at length scales $\gg R_\psi$.

\textbf{Conjecture:} Under these assumptions, the UBT field equations
enforce chronology protection at least as strongly as in standard GR.
The Cauchy-horizon divergence is UV-regulated by the $\psi$ sector but
remains non-integrable, destroying the horizon before CTC formation.
\end{conjecture}

\subsection{Scenario B: Circumvention of Chronology Protection}
If the imaginary sector $\psi$ acts as a \emph{channel} that allows
energy conditions to be violated locally (Assumption 4 fails) while the
Cauchy-horizon stress-energy remains finite (Assumption 3 modified),
then the standard back-reaction argument fails and CTCs may form.

\begin{conjecture}[Chronology protection circumvention]
\label{conj:circumvention}
\textbf{Assumptions:}
(a) $\psi$ is a second independent time direction (Interpretation A);
(b) the imaginary metric sector $h_{\mu\nu}$ provides negative energy
density at the Cauchy horizon;
(c) the $\psi$-sector regulates the divergence of
$\langle T_{\mu\nu}\rangle_{\mathrm{ren}}$ to a finite value;
(d) back-reaction of this finite stress-energy is insufficient to
destroy the horizon.

\textbf{Conjecture:} Under these assumptions, the UBT field equations
do \emph{not} enforce chronology protection in the imaginary sector,
and closed timelike curves in the $(t,\psi)$ plane are not forbidden.
\end{conjecture}

\begin{remark}
Conjectures~\ref{conj:enhanced_cp} and \ref{conj:circumvention} are
mutually exclusive: they correspond to different interpretations of $\psi$
(Interpretation B versus A, respectively).  Determining which interpretation
is physically realised requires solving the full biquaternionic field
equations in a CTC-enabling background, which remains an open problem.
\end{remark}

%----------------------------------------------------------------------
\section{Role of the $\psi$-Sector as Regulator}
%----------------------------------------------------------------------
The compact $\psi$ direction with radius $R_\psi$ introduces a
Kaluza–Klein-like tower of modes with masses
\begin{equation}
  m_n^{(\psi)} \;=\; \frac{n}{R_\psi}, \qquad n \in \mathbb{Z}.
  \label{eq:kk_masses}
\end{equation}
These massive modes suppress quantum fluctuations in the imaginary direction
at energies below $R_\psi^{-1}$.  The renormalised stress-energy on the
Cauchy horizon receives contributions from all KK modes; whether these
sum to a divergence or remain finite depends on the spectral properties
of the UBT propagator in the biquaternionic background.

A key open question is whether the KK sum
\begin{equation}
  \langle T_{\mu\nu}\rangle_{\mathrm{ren}}^{\psi} \;=\;
  \sum_{n=-\infty}^{\infty}
  \langle T_{\mu\nu}\rangle_n^{\mathrm{Cauchy}}
  \label{eq:kk_sum}
\end{equation}
converges or diverges.  If convergent, the back-reaction is finite and
Scenario~B (circumvention) is plausible.  If divergent, the back-reaction
is enhanced and Scenario~A (strengthened chronology protection) applies.

%----------------------------------------------------------------------
\section{Summary and Open Questions}
%----------------------------------------------------------------------
\begin{itemize}
  \item The standard QFT assumptions underlying Hawking's chronology
    protection conjecture are \emph{not all satisfied} in UBT.
  \item Under Interpretation B ($\psi$ = internal phase), chronology
    protection is at least as strong as in GR (Conjecture~\ref{conj:enhanced_cp}).
  \item Under Interpretation A ($\psi$ = second time direction),
    chronology protection may be circumvented if the imaginary sector
    provides a finite regulator for the Cauchy-horizon stress-energy
    (Conjecture~\ref{conj:circumvention}).
  \item The key open question is whether the KK sum \eqref{eq:kk_sum}
    converges in the UBT propagator.  This requires explicit calculation
    in a biquaternionic CTC background.
\end{itemize}

\medskip
\noindent\textbf{All statements in this document are speculative.  No
claim constitutes an established result of the canonical UBT framework.
The conjectures above are contingent on listed assumptions and must not
be promoted to theorems without independent proof.}

%----------------------------------------------------------------------
\bibliographystyle{plain}
\begin{thebibliography}{9}
\bibitem{Hawking1992}
  S.~W.~Hawking,
  \emph{Chronology protection conjecture},
  Phys.\ Rev.\ D \textbf{46}, 603 (1992).
\bibitem{Kay1997}
  B.~S.~Kay, M.~J.~Radzikowski and R.~M.~Wald,
  \emph{Quantum field theory on spacetimes with a compactly generated
  Cauchy horizon},
  Commun.\ Math.\ Phys.\ \textbf{183}, 533 (1997).
\bibitem{Visser1995}
  M.~Visser,
  \emph{Lorentzian Wormholes: from Einstein to Hawking},
  AIP Press (1995).
\bibitem{UBT_core}
  D.~Jaroš,
  \emph{Unified Biquaternion Theory: core field equations},
  Repository document \texttt{consolidation\_project/ubt\_core\_main.tex} (2025).
\bibitem{CTC_ST3}
  D.~Jaroš,
  \emph{ST-3: CTC Conditions in the Biquaternionic Metric},
  Repository document
  \texttt{speculative\_extensions/causality/ctc\_conditions.tex} (2026).
\end{thebibliography}

\end{document}
